\subsection{The constructive models and their exact comparison limits}
\paragraph{\fid{F18}: weight lattices and rational semigroups.}
Regard $g(\nu,(d,d),(d,d))$ as the coefficient of
$\det^d\boxtimes\det^d$ in the restriction of $S_\nu(V\otimes W)$,
where $V,W$ have rank two. This is an equality in the representation
$K_0$ of the group scheme $\GL(V)\times\GL(W)$ after passage to
characteristic zero. Characters of a diagonalizable group correspond
to gradings (SGA~3, I, 4.7.3). Integral difference-constraint arrays
would supply a particular saturated graded semigroup. The flooring
and dilation proof of \fid{F18} transfers unchanged to that weight
lattice. Its parity example requires rational-lattice or congruence
data instead. SGA~6's formal $\lambda$ identities do not supply those
data; cote~133, scans~18--21, is a vocabulary anchor rather than a
positive-rule theorem.

\paragraph{\fid{F19}, \fid{F21}, \fid{F25}: coherent algebra and its actual module.}
Here there is a specified scheme over a localization of
$\Spec\Z[v,v^{-1}]$: the finite flat center $S$ of \fid{F19}, with
the explicitly trivial bundle $E$ of ranks $1,1,2$.
The noncommutative algebra is $\Gamma(S,\End E)$, and the actual
right tensor module is $\Gamma(S,\mathcal Hom(E,\mathcal F))$ for
the displayed coherent branch sheaf $\mathcal F$. Thus the module
dictionary is an explicit action-preserving isomorphism, not just
an equality of characters. The matrices in \fid{F21} carry the proof
over to this sheaf language. Equivariant quasi-coherent modules are
the framework of cote~90, scan~17, and SGA~3, I, 4.7.2.

The branch modules $L_i$ are coherent but are not asserted to be
locally free or perfect over the singular center. Their natural
Grothendieck group is $G_0(C)$; one must not silently apply vector-bundle
$\lambda$ operations from $K_0(C)$ to them. The periodic resolution
in \fid{F25} computes their actual coherent derived Hom. Its
$R/(\Delta)$ records an order-two deformation invisible in the
constant coefficient realization. The regular algebra basis in that
finding is unrelated to any assertion that this singular coherent
category is semisimple. Its scalar $\lambda$ has compact-support motive
\[
 \mathbb L^{-2}(\mathbb L-1)^2(1+\mathbb L+\mathbb L^2)
\]
in the virtual Tate ring. The trace equality follows from
$\mathbb G_m^2\times\PP^2$ and the trace formula; its mixed signs are
cohomological Euler signs, not dimensions of one pure degree.

\paragraph{\fid{F22}: the \'etale topos and the structure sheaf retain different data.}
The finite maps $\Spec C_n\to\Spec C_m$ are universal
homeomorphisms. EGA~IV, 18.1.2, in the precise modern formulation
Stacks 04DZ/03SI, transports the entire ordinary coefficient-sheaf
topos over the trait. The proof in \fid{F22} explicitly identifies
the constant nearby-cycle triangle with
$\Q_\ell\to\Q_\ell^2\to\Q_\ell$; all three objects have scalar
weight-zero Frobenius. In contrast the conductor and crossed coherent
$\operatorname{Ext}^1$ retain $m$. This is an exact instance of the
difference between a coefficient topos and a ringed coherent space.
It does not forbid adding a separately specified nilpotent action to
a Weil sheaf. No manuscript statement is claimed to prove this
nonstandard deformation comparison.

\paragraph{\fid{F20}: effective Schur operations and Tate coefficient motives.}
On a split affine Grassmannian, the primitive Young projector is an
actual idempotent in a rational symmetric tensor category. Its image
is an effective object, and its class is the specified Schur operation
on the locally defined representation. The modern motivic Satake
theorems cited in \fid{F20} identify its coefficient motive as
$\mathbf1(-h)^{\oplus a}$, not merely an alternating virtual class.
The ordinary coefficient realization has scalar Frobenius $q^h$.
This supplies an actual motive in these classical cases; it does not
assert the existence of a Chow motive for every singular orbit closure.
The source of the motivic theorem is Richarz--Scholbach, not SGA~6
or cote~16. The latter's scans~35,37,39 supply the historical language
of algebraic projectors. Canonical bases and efficient rank labels
remain separate information even when the motive is Tate.

\paragraph{\fid{F23}: induction in the effective representation monoid.}
A signed orbit is the representation of a constant finite group
scheme induced from a one-dimensional stabilizer representation.
Over $\Z[1/d!]$ with the relevant integral lattices, finite-group
averaging splits invariant projectors; over characteristic zero
the representation category is semisimple. The direct-sum equality
of \fid{F23} is effective in representation $K_0$, and its Young
stabilizer specialization is the ordinary symmetric/exterior Pieri
operation. The proof by orbit modules and restriction of Schur
functors is unchanged in this language. The pure Tate $\GL_2$
example shows why the $S_d$ correspondence action is additional to
the weight and motive. Cotes~90/17 and~133/19--21 give the respective
module and finite-set-operation anchors; neither provides the missing
universal signed basis.

\paragraph{\fid{F24}: a signed class can still retain a largest eigenvalue.}
For a bounded Weil complex, group equal roots in its alternating trace.
Its visible multiplicities are integers, possibly of both signs.
The external tensor product maps this class to its power in the
group ring of eigenvalues. The augmentation-and-Vandermonde argument
in \fid{F24} proves that an even power retains a positive top-radius
coefficient. This transfers the proof exactly to the triangulated
$K_0$ trace formalism of SGA~4/5, without requiring one input degree.
The negative visible path-output spectrum and any reciprocal visible
duality are still hypotheses. A zero virtual class may erase actual
cohomology, and finite extensions can erase a signed top spectrum;
the two countermodels specify those limits.

\paragraph{\fid{F26}--\fid{F27}: localization, bounded complexes, and finite traces.}
Finite Weil cohomology gives rational resolvents and zeta factors
$\det(1-TF)$, including their alternating products. This is the
finite determinant mechanism underlying the trace formalism; the
non-Laurent rational grading function in \fid{F26} cannot be realized
by such a fixed finite trace in all extensions. Its recurrence proof
is independent of any particular choice of category. The generic
Fredholm determinant formulas visible in cote~1, scans~77--78,
do not identify it with a finite Weil object or a $\Xi$-function.

Bounded complexes in the Hecke category do realize the signed Euler
class $P$. Their four basis coordinates prove the one-parity
obstruction in \fid{F27}. The rank-one bimodule comes from the
finite flat correspondence
$\Spec A\times_{\Spec A^s}\Spec A$; passing to its normalization
separates the character points and yields the other positive cone.
The actual four-point Tate construction retains the rank-one action
after the stated localization. Algebraic localization, coherent
normalization, and a trace-compatible tensor realization are three
specified operations here; none is inferred from the others.

